\documentclass[paper=screen,mode=present,style=zysimple]{powerdot}

\usepackage{fancyvrb,verbatim,listings}
\usepackage{proof}
\usepackage[all]{xypic,xy}

\usepackage{latexsym,stmaryrd,amsmath,amsfonts}

%% \newtheorem{definition}{Definition}
%% \newtheorem{theorem}{Theorem}
%% \newtheorem{lemma}{Lemma}
\newtheorem{prop}{Proposition}
%% \newtheorem{corollary}{Corollary}
%% \newenvironment{proof}{{\sc Proof . }}{{\hfill $\square$}}
%% \newcommand{\myqed}{\hfill $\square$}
%% \newcommand{\ped}{\vspace{-1.6em}}
%% \newcommand{\note}[1]{\noindent\fbox{\parbox{\textwidth}{\em #1}}}

\newcommand{\abind}[2]{\mathtt{bind} \ (\, {#1} \,)  \ \mathtt{in} \ {#2}}
\newcommand{\acase}{\mathtt{case}}
\newcommand{\adv}{\mathcal{A}}
\newcommand{\aspecto}[1]{\xrightarrow{\; {#1} \;}}
\newcommand{\asplinear}{\mathfrak{l}}
\newcommand{\aspmodal}{\mathfrak{m}}
\newcommand{\aspnormal}{\mathfrak{n}}
%
\newcommand{\betareduce}{\leadsto_\beta}
\newcommand{\betareduceto}{\twoheadrightarrow_\beta}
\newcommand{\bit}{\mathtt{B}}
\newcommand{\bitone}{\mathtt{B}_1}
\newcommand{\bitzero}{\mathtt{B}_0}
%
\newcommand{\caseto}{\triangleright}
\newcommand{\cbind}[3]{\mathtt{bind} \ {#1} = {#2} \ \mathtt{in} \ {#3}}
\newcommand{\ccase}[4]{\mathtt{case} \; {#1} \; \mathtt{of} \; \cnil \caseto {#2}; \;
                                          \bitzero x_0 \caseto {#3}; \; \bitone x_1 \caseto {#4}}
\newcommand{\cequiv}{\simeq}
\newcommand{\cif}[3]{\mathtt{if} \; {#1} \; \mathtt{then} \; {#2} \; \mathtt{else} \; {#3}}
\newcommand{\ctrue}{\mathtt{true}}
\newcommand{\cfalse}{\mathtt{false}}
\newcommand{\cfix}{\mathtt{fix}\,}
\newcommand{\cfixin}[2]{\mathtt{fix}\; {#1}\; \mathtt{in} \; {#2}}
\newcommand{\cfold}{\mathtt{fold}\,}
\newcommand{\chd}{\mathtt{hd}}
\newcommand{\cinj}{\mathtt{inj}}
\newcommand{\clift}[1]{\lfloor{#1}\rfloor}
\newcommand{\clet}[3]{\mathtt{let} \ {#1} = {#2} \ \mathtt{in} \ {#3}}
\newcommand{\cletst}[3]{\mathtt{liftlet} \ {#1} = {#2} \ \mathtt{in} \ {#3}}
\newcommand{\cletrec}[3]{\mathtt{letrec} \ {#1} = {#2} \ \mathtt{in} \ {#3}}
\newcommand{\cnil}{\mathtt{nil}}
\newcommand{\conc}{\! \bullet \!}
\newcommand{\concat}{\pmb{\mathit{conc}}}
\newcommand{\context}{\mathbb{C}}
\newcommand{\cop}{\;\mathtt{op}\;}
\newcommand{\cpair}[1]{\langle{#1}\rangle}
\newcommand{\cproj}{\mathtt{proj}}
\newcommand{\crand}{\mathtt{rand}}
\newcommand{\crec}{\mathtt{rec}}
\newcommand{\croll}{\mathtt{roll}\,}
\newcommand{\crule}[1]{\mbox{{\textsl{#1}}}}
\newcommand{\cslrtype}{\Vdash}
\newcommand{\csucc}{\mathtt{S}}
\newcommand{\ctl}{\mathtt{tl}}
\newcommand{\cut}{\pmb{\mathit{cut}}}
\newcommand{\cval}{\mathtt{return}}
\newcommand{\cwhile}[2]{\mathtt{while}\; {#1} \; \mathtt{do} \; {#2} \; \mathtt{od}}
\newcommand{\cvalof}{\pmb{valueof}}
\newcommand{\cunfold}{\mathtt{unfold}\,}
\newcommand{\cunit}{\mathtt{unit}}
\newcommand{\cunroll}{\mathtt{unroll}\,}
\newcommand{\cupdate}{\pmb{update}}
\newcommand{\czero}{\mathtt{O}}
%
\newcommand{\deduces}{\vdash}
\newcommand{\defeq}{\stackrel{\mathrm{def}}{=}}
\newcommand{\delay}[1]{\left\lfloor{#1}\right\rfloor}
\newcommand{\distr}{\mathcal{D}}
%
\newcommand{\eagerreduce}{\xrightarrow{\mathrm{eager}}}
\newcommand{\Eagerreduce}{\stackrel{\mathrm{eager}}{\Reduce}}
\newcommand{\eqtest}{\stackrel{\mbox{\tiny ?}}{=}}
\newcommand{\Eval}[1]{{\left\llbracket{#1}\right\rrbracket}}
\newcommand{\evalcontext}{\mathbb{E}}
\newcommand{\Example}{\emgreen{\pmb{\textsc{Example}\;}}}
\newcommand{\expansion}[1]{\tilde{{#1}}}
%
\newcommand{\ford}{\mathrm{ord}}
%
\newcommand{\getval}{\leftarrow}
%
\newcommand{\head}{\pmb{\mathit{hd}}}
\newcommand{\termhole}{[\;]}
\newcommand{\hoaretriple}[3]{\{{#1}\} \; {#2} \; \{{#3}\}}
%
\newcommand{\imply}{\Rightarrow}
\newcommand{\indist}{\simeq}
\newcommand{\iteration}{\pmb{\mathit{it}}}
%
\newcommand{\lazyreduce}{\xrightarrow{\mathrm{lazy}}}
\newcommand{\Lazyreduce}{\lazyreduce\!\!\!\!\!\to}
\newcommand{\ldot}{\, . \, }
\newcommand{\linearto}{\multimap}
\newcommand{\logrel}{(\mathcal{R}_\tau)_{\mbox{$\tau$ type}}}
\newcommand{\lr}[1]{\mathcal{R}_{#1}}
\newcommand{\lub}{\sqcup}
%
\newcommand{\metalambda}{\underline{\lambda}}
\newcommand{\modal}{\square}
%
\newcommand{\nat}{\mathbb{N}}
\newcommand{\nReduce}{\nrightarrow\!\!\!\!\to}
\newcommand{\nreduceto}{\nrightsquigarrow}
\newcommand{\nullenv}{\varnothing}
%
\newcommand{\orc}{\mathcal{O}}
\newcommand{\oslrtype}{\Vdash}
%
\newcommand{\porder}{\sqsubseteq}
\newcommand{\prefix}{\pmb{\mathit{pref}}}
\newcommand{\proba}{\mathbf{Pr}}
%
\newcommand{\randlen}{\pmb{\mathit{rl}}}
\newcommand{\randnum}{\pmb{\mathit{rn}}}
\newcommand{\randstr}{\pmb{\mathit{rs}}}
\newcommand{\reduce}{\leadsto}
\newcommand{\reduceto}{\rightsquigarrow\hspace*{-6px}\shortrightarrow}
\newcommand{\Reduce}{\to\!\!\!\to}
\newcommand{\Reduceto}{\Downarrow}
%
\newcommand{\set}[1]{\{{#1}\}}
\newcommand{\setbits}{\mathbb{B}}
\newcommand{\slrtype}{\vdash}
\newcommand{\splt}{\pmb{\mathit{split}}}
\newcommand{\suffix}{\pmb{\mathit{suff}}}
\newcommand{\sysname}{\pmb{\mathcal{C}}}
%
\newcommand{\tail}{\pmb{\mathit{tl}}}
\newcommand{\tbits}{\mathsf{Bits}}
\newcommand{\tbool}{\mathsf{Bool}}
\newcommand{\tnat}{\mathsf{Nat}}
\newcommand{\tcomp}{\mathsf{T}}
\newcommand{\tlift}[1]{\lfloor{#1}\rfloor}
\newcommand{\tlist}{\mathsf{List}}
\newcommand{\tself}{\mathsf{self}}
\newcommand{\tstore}{\mathsf{Store}}
\newcommand{\tstream}{\mathsf{Stream}}
\newcommand{\tloc}{\mathsf{Loc}}
\newcommand{\ttcase}{\mathtt{case}}
\newcommand{\ttif}{\mathtt{if}}
\newcommand{\ttthen}{\mathtt{then}}
\newcommand{\ttelse}{\mathtt{else}}
\newcommand{\ttin}{\mathtt{in}}
\newcommand{\ttlet}{\mathtt{let}}
\newcommand{\ttof}{\mathtt{of}}
\newcommand{\tunit}{\mathsf{Unit}}
\newcommand{\Type}{\mathcal{T}\!\mathit{ype}}
%
\newcommand{\vbool}{\pmb{\mathbb{B}}}
\newcommand{\vff}{\mathfrak{f}\!\mathfrak{f}}
\newcommand{\vnat}{\pmb{\mathrm{N}}}
\newcommand{\vnil}{\epsilon}
\newcommand{\vtt}{\mathfrak{t}\!\mathfrak{t}}
%
\newcommand{\weakpre}{\pmb{wp}}




\pdsetup{lf={Theory and Practice of Functional Programming}}

\title{{\small \textsc{Theory and Practice of Functional Programming}}
\\[12pt]
{\Large \color{TitleColor} Topics in Type Theory}
}
\author{\small \textbf{Dr. ZHANG Yu}
  \\[2pt] \small Institute of Software, Chinese Academy of Sciences 
  \\[5pt] \includegraphics[height=.04\slideheight]{iscas.eps} 
  \qquad \qquad 
  \includegraphics[height=.04\slideheight]{cas.eps}
}
\date{\scriptsize Fall term, 2011 \qquad GUCAS, Beijing}

\newcommand{\RED}{\mathsf{RED}}

\begin{document}

\maketitle 

\section{Type Inference} 

\begin{slide}[method=direct]{Principal types}
\begin{itemize}
\item A few type judgements:
\begin{smalltext}
\begin{eqnarray*} 
&& 
\vdash \lambda x : \tnat \ldot \lambda y : \tnat \to \tnat \ldot y x : 
\emblue{\tnat} \to (\emgreen{\tnat \to \tnat}) \to \empurple{\tnat}
\\ && 
\vdash \lambda x : \tbool \ldot \lambda y : \tbool \to \tnat \ldot y x : 
\emblue{\tbool} \to (\emgreen{\tbool \to \tnat}) \to \empurple{\tnat}
\\ && 
\vdash \lambda x : (\tnat \to \tbool) \ldot \lambda y : (\tnat \to \tbool) \to \tbool \to \tbool \ldot y x : 
\\ && \qquad
(\emblue{\tnat \to \tbool}) \to (\emgreen{(\tnat \to \tbool) \to \tbool \to \tbool}) \to \empurple{\tbool \to \tbool}
\end{eqnarray*}
\end{smalltext}
All types above are of the pattern $\emblue{X} \to (\emgreen{X \to Y}) \to \empurple{Y}$.
\item In Haskell, the system will give the following type to the term $\lambda x \ldot \lambda y \ldot y x$:
  %% # fun x -> fun y -> y x ;;
  %% - : 'a -> ('a -> 'b) -> 'b = <fun>
\vspace*{-1.8em}
\begin{smalltext}
\begin{verbatim}
    Prelude> :t \x -> \y -> y x
    \x -> \y -> y x :: t1 -> (t1 -> t) -> t
\end{verbatim}
\end{smalltext}
\item $X \to (X \to Y) \to Y$ is the most general type (or the \emblue{principal type}) for the term 
$\lambda x \ldot \lambda y \ldot y x$.
\\
Every typable lambda term in the simply typed lambda calculus has a principal type.
\end{itemize}
\end{slide}

\begin{slide}{Type templates}
\begin{itemize}
\item A \emblue{type template} is a simple type, built from \emblue{type variables} and possibly base types.
\[
T, S, \ldots ::= X \mid b \mid T \to S \mid T \times S \mid T + S
\]
where $X$ is type variable, $b$ ranges over a set of type constants.
\item A \emblue{type substitution} is a function from type variables to type templates: 
\vspace*{-0.5em}
  $$[X_1 \mapsto T_1, \ldots, X_n \mapsto T_n]  \vspace*{-0.5em}$$
We define the template $\sigma(T)$, obtained by applying the substitution $\sigma$ to templpate $T$, 
as follows:
\begin{smalltext}
\vspace*{-0.5em}
\begin{eqnarray*}
\sigma(X) & = & T, \mbox{ if $X \mapsto T \in \sigma$}
\\
\sigma (b) & = & b
\\
\sigma (T \to S) & = & \sigma(T) \to \sigma(S) 
\\
\sigma (T \times S) & = & \sigma(T) \times \sigma(S) 
\\[-2em]
\end{eqnarray*}
\end{smalltext}
A template $T$ is \emblue{more general} than $S$ if there exists a substitution $\sigma$ 
such that $\sigma(T) = S$, or $S$ instantiates $T$.
\end{itemize}
\end{slide}

\begin{slide}{Unification}
\begin{itemize}
\item Given an equation between type templates, a \emblue{unifier} is a type substitution
  that solves the equation. 
\begin{itemize}\small
\item $X \to (X \to Y) = (Y \to Z) \to W$ has a solution 
\[
\sigma = [ X \mapsto Y \to Z, W \mapsto (Y \to Z) \to Y ]
\vspace*{-1.5em}
\]
\item The equation $X \times (X \to Z) = (Z \to Y) \times Y$ does not have any solution.
\end{itemize}
\item A unifier of two sequences of type templates $\tilde{T} = T_1, \ldots , T_n$ and 
  $\tilde S = S_1, \ldots , S_n$, is a type substitution $\sigma$ s.t. 
  $\sigma(T_i) = \sigma(S_i)$ for all $i$.
\item A types substitution is a \emblue{most general unifier} of $T = S$ if it is a unifier 
and it is more general than any other unifier of $T =S$. \\
\begin{itemize} \small
\item A type substitution $\sigma$ is {\em more general} than $\sigma'$ if there exists another 
substitution $\rho$ such that $\rho \circ \sigma = \sigma'$.
\item Let $\sigma = [X_1 \mapsto T_1, \ldots , X_n \mapsto T_n]$. 
\emblue{$$\rho \circ \sigma = [X_1 \mapsto \rho(T_1), \ldots, X_n \mapsto \rho (T_n)] $$}
\end{itemize}
\end{itemize}
\end{slide}

\begin{slide}{The unification algorithm}
\begin{itemize}
\item \emblue{Unification} is the process of determining a most general unifier.
\vspace*{-0.5em}
\begin{itemize}\small
\item Input: two sequences of type templates $\tilde T$ and $\tilde S$.
\item Output: a most general unifier $\sigma$ if it exists.
\vspace*{-1em}
\end{itemize}
\item The most general unification algorithm \emred{$MGU$}:
\vspace*{-0.5em}
\begin{itemize}\small
\item $MGU(T; T) = \mathrm{id}$, if $T$ is a variable or a base type;
\item $MGU(X; S) = [ X \mapsto S ]$, if \empink{$X$ does not occur in $S$};
\item $MGU(T; Y) = [ Y \mapsto T ]$, if \empink{$Y$ does not occur in $T$};
\item $MGU(T_1 \to T_2; S_1 \to S_2) = MGU (T_1, T_2; S_1, S_2)$;
\item $MGU(T_1 \times T_2; S_1 \times S_2) = MGU (T_1, T_2; S_1, S_2)$;
\item $MGU(T_1 + T_2; S_1 + S_2) = MGU (T_1, T_2; S_1, S_2)$;
\item $MGU(T; S)$ fails in all other cases;
\item $MGU(T, \tilde T; S, \tilde S) = \sigma' \circ \sigma$, where $\sigma = MGU(\tilde T; \tilde S)$ 
  and $\sigma' = MGU(\sigma(T); \sigma(S))$.
\vspace*{-0.5em}
\end{itemize}
\item If $MGU(\tilde T; \tilde S) = \sigma$, then $\sigma$ is a most general unifier of $\tilde T$ 
  and $\tilde S$.
\end{itemize}
\end{slide}

\begin{slide}{The type inference algorithm}
The type inference algorithm $TI$: \emblue{$TI(\Gamma \vdash e: {T})$} is a most general substitution $\sigma$ 
such that \emred{$\sigma(\Gamma) \vdash e : \sigma(T)$ holds} ($TI$ fails if no such $\sigma$).
\vspace*{-0.5em}
\begin{itemize}\small
\item[--] $TI(x_1 : T_1, \ldots , x_n : T_n \vdash x_i : T) = MGU(T_i; T)$. \vspace*{-0.5em}
\item[--] $TI(\Gamma \vdash e_1 + e_2 : T) = \sigma_2 \circ \sigma_1 \circ \sigma$, where \\
  \begin{minipage}{.5\textwidth}
  \begin{itemize}
  \item $\sigma = MGU(\tnat; T)$
  \item $\sigma_1 = TI(\sigma(\Gamma) \vdash e_1 : \tnat)$
  \item $\sigma_2 = TI(\sigma_1 \circ \sigma(\Gamma) \vdash e_1 : \tnat)$
  \end{itemize}
  \end{minipage}
  \hfill
  \begin{minipage}{.3\textwidth}
    $$\prooftree \Gamma \vdash e_1 : \tnat \quad \Gamma \vdash e_2: \tnat \justifies 
    \Gamma \vdash e_1 + e_2 : \tnat \endprooftree$$
  \end{minipage} 
  \vspace*{-1em}
\item[--] $TI(\Gamma \vdash e_1 e_2 : T) = \sigma_2 \circ \sigma_1$, where \\
  \begin{minipage}{.5\textwidth}
  \begin{itemize}
  \item $\sigma_1 = TI(\Gamma \vdash e_1 : X \to T)$
  \item $\sigma_2 = TI(\emblue{\sigma_1(\Gamma)} \vdash e_2 : \emblue{\sigma_1(X)})$ 
  \item $X$ is a fresh type variable
  \end{itemize}
  \end{minipage}
  \hfill
  \begin{minipage}{.3\textwidth}
    $$\prooftree \Gamma \vdash e_1 : \tau \to \tau' \quad \Gamma \vdash e_2: \tau \justifies 
    \Gamma \vdash e_1 \ e_2 : \tau' \endprooftree$$ 
  \end{minipage}
  \vspace*{-1em}
\item[--] $TI(\Gamma \vdash \lambda x: S \ldot e : T) = \sigma_2 \circ \sigma_1$,  where \\
  \begin{minipage}{.6\textwidth}
  \begin{itemize}
  \item $\sigma_1 = MGU(T; S \to X)$
  \item $\sigma_2 = TI(\sigma_1(\Gamma), x: \sigma_1(S) \vdash e : \sigma_1(X))$
  \item $X$ is a fresh variable
  \end{itemize}
  \end{minipage}
  \hfill
  \begin{minipage}{.3\textwidth}
    $$\prooftree \Gamma, x : \tau \vdash e : \tau' 
    \justifies \Gamma \vdash \lambda x \ldot e : \tau \to \tau' \endprooftree$$
  \end{minipage}
\end{itemize}
\end{slide}

\begin{slide}{The type inference algorithm}
\begin{itemize}\small
\item[--] $TI(\Gamma \vdash \cpair{e_1, e_2} : T) = \sigma_2 \circ \sigma_1 \circ \sigma$, where 
  \begin{minipage}{.5\textwidth}
  \begin{itemize}
  \item $\sigma = MGU(T; X \times Y)$
  \item $\sigma_1 = TI(\sigma(\Gamma) \vdash e_1 : \sigma(X))$
  \item $\sigma_2 = TI(\sigma_1 \circ \sigma (\Gamma) \vdash e_2 : \sigma_1 \circ \sigma(Y))$ 
  \item  $X, Y$ are fresh type variables
  \end{itemize}
  \end{minipage}
  \begin{minipage}{.3\textwidth}
    $$\prooftree \Gamma \vdash e_1 : \tau_1 \quad \Gamma \vdash e_2 : \tau_2
    \justifies \Gamma \vdash \cpair{e_1, e_2} : \tau_1 \times \tau_2 \endprooftree$$
  \end{minipage}
  \vspace*{-1em}
\item[--] $TI(\Gamma \vdash \cproj_1 e : T) = TI(\Gamma \vdash e : T \times X)$, where 
  \begin{minipage}{.5\textwidth}
  \begin{itemize}
  \item $X$ is a fresh type variable
  \item $TI(\Gamma \vdash \cproj_2 e : T)$ is similar.
  \end{itemize}
  \end{minipage}
  \begin{minipage}{.3\textwidth}
    $$\prooftree \Gamma \vdash e : \tau_1 \times \tau_2
    \justifies \Gamma \vdash \cproj_i \ e : \tau_i \endprooftree$$
  \end{minipage}
  \vspace*{-1em}
\item[--] $TI(\Gamma \vdash \cinj_1 e : T) = TI(\Gamma \vdash e : T \times X)$, where 
  \begin{minipage}{.5\textwidth}
  \begin{itemize}
  \item $\sigma = MGU(X + Y; T)$ 
  \item $\sigma_1 = TI(\sigma{\Gamma} \vdash e : X$
  \item $X, Y$ is a fresh type variable
  \item $TI(\Gamma \vdash \cinj_2 e : T)$ is similar.
  \end{itemize}
  \end{minipage}
  \begin{minipage}{.3\textwidth}
    $$\prooftree \Gamma \vdash e : \tau_i
    \justifies \Gamma \vdash \cinj_i \ e : \tau_1 + \tau_2 \endprooftree$$
  \end{minipage}
\end{itemize}
\end{slide}

\begin{slide}{The type inference algorithm}
\begin{itemize}\small
\item[--] $TI(\Gamma \vdash \mathtt{case} \; e \; \mathtt{of} \; \cinj_1 x_1 \; \mathtt{in} \; e_1 ; \cinj_2 x_2 \; \mathtt{in} \; e_2 : T) = \sigma_2 \circ \sigma_1 \circ \sigma$, where 
  \begin{itemize}
  \item $\sigma = TI(\Gamma \vdash e :X + Y)$
  \item $\sigma_1 = TI(\sigma(\Gamma), x_1 : \sigma(X) \vdash e_1 : \sigma(T))$
  \item $\sigma_2 = TI(\sigma_1 \circ \sigma (\Gamma), x_2 : \sigma_1 \circ \sigma (Y) \vdash e_2 : \sigma_1 \circ \sigma(Y))$ 
  \item  $X, Y$ are fresh type variables
  \end{itemize}
$$
\prooftree
\Gamma, x_1 : \tau_1 \vdash e_1 : \tau 
\quad 
\Gamma, x_2 : \tau_2 \vdash e_2 : \tau 
\quad 
\Gamma \vdash e : \tau_1 + \tau_2
\justifies
\Gamma \vdash 
\mathtt{case} \; e \; \mathtt{of} \; \cinj_1 x_1 \; \mathtt{in} \; e_1 ; \cinj_2 x_2 \; \mathtt{in} \; e_2 : \tau
\endprooftree
$$
\item[--] $TI(\Gamma \vdash \cfix e : T) = TI(\Gamma \vdash e : T \to T)$.
$$
  \prooftree \Gamma \vdash e : \tau \to \tau 
  \justifies \Gamma \vdash \cfix \ e : \tau \endprooftree
$$
\end{itemize}
%% \item $TI(\Gamma \vdash e : T)$ returns a most general substitution $\sigma$ such that
%%   $\sigma(\Gamma) \vdash e : \sigma(T)$ holds. It fails if no such substitution exists.
In real applications, type inference only takes two arguments $\Gamma$ and $e$, 
and we do inference by $TI(\Gamma \vdash e : \emred{X})$, with $X$ a fresh type variable.
\end{slide}

\section{Polymorphism and System F}

\begin{slide}{Polymorphism and types}
\begin{itemize}
\item Programs with generic behaviours that are independent of types:
\begin{eqnarray*}
&& {id}_{\tau} \defeq \lambda x : \tau \ldot x
\\ && 
{compose}_{\tau_1, \tau_2, \tau_3} \defeq 
\lambda f : \tau_2 \to \tau_3 \ldot \lambda g : \tau_1 \to \tau_2 \ldot \lambda x : \tau_1 \ldot f (g (x))
\\ & & 
\ldots \quad \ldots
\\[-3em]
\end{eqnarray*}
\item These expressions can be considered as polymorphic, and can be given a unifying type in a polymorphic type system.
\end{itemize}
\end{slide}

\begin{slide}{System F}
\begin{itemize}
\item System F \emgreen{[Jean-Yves Girard]}.
\begin{itemize}
\item Extension of simply-typed $\lambda$-calculus with polymorphic types.
\item A.k.a. {\em polymorphic $\lambda$-calculus}, independently by \emgreen{John Reynolds}.
\end{itemize}
\item Types
\[
\begin{array}{lrl@{\quad}l}
\tau, \tau', \ldots & ::= & X \mid \tau \to \tau' & 
\\ 
& \mid & \emred{\forall X \ldot \tau} & \mbox{polymorphic type}
\end{array}
\]
\item Terms (expressions)
\[
\begin{array}{lrl@{\quad}l}
e, e' & ::= & x \mid \lambda x : \tau \ldot e \mid e \, e' 
\\
& \mid & \emred{\Lambda X \ldot e} & \mbox{universal abstraction}
\\
& \mid & \emred{e [\tau]} & \mbox{universal application}
\end{array}
\]
\end{itemize}
\end{slide}

\begin{slide}{System F}
\begin{itemize}
\item Universal reduction 
\[
\emblue{(\Lambda X \ldot e) \tau \reduce e [\tau / X]}
\vspace*{-2em}
\]
\item Typing
\begin{itemize}
\item Type information judgement (well-formed type)
\[
\prooftree
\Delta , X : \Type \vdash \tau : \Type
\justifies
\Delta \vdash \forall X \ldot \tau : \Type
\endprooftree
\]
\item Typing judgement
\[
\begin{array}{c}
\prooftree
\Delta, X : \Type ; \Gamma \vdash e : \tau
\justifies
\Delta ; \Gamma \vdash \Lambda X \ldot e : \forall X \ldot \tau
\endprooftree
\qquad
\prooftree
\Delta ; \Gamma \vdash e : \forall X \ldot \tau : \Type
\quad 
\Delta \vdash \tau' : \Type
\justifies
\Delta ; \Gamma \vdash e [\tau'] : \tau[\tau' / X]
\endprooftree
\end{array}
\]
\end{itemize}
\end{itemize}
\end{slide}

\begin{slide}{Representing simple types}
\begin{itemize}
\item Unit type 
\[
\emblue{\tunit \defeq \forall X \ldot X \to X}
\]
This type contains one and \emred{only one} element --- the identity function:
\[
* \defeq \Lambda X \ldot \lambda x : X \ldot x
\]
This is supported by the \emblue{parametricity theorem}, which says that every program of type 
$\forall X \ldot X \to X$ is {\em observationally equivalent} to the above identity function.
\item Empty type
\[
\emblue{\bot \defeq \forall X \ldot X}
\]
There is no well-typed program of this type. 
\\
Polymorphism says that a program of the type returns a value, independently of the choice of type $X$, 
including types with no programs.
\end{itemize}
\end{slide}

\begin{slide}{Representing simple types}
\begin{itemize}
\item Boolean 
\[
\emblue{\tbool \defeq \forall X \ldot X \to X \to X}
\]
The only two elements of $\tbool$ 
\[
\ctrue \defeq \Lambda X \ldot \lambda x : X \ldot \lambda y : X \ldot x , 
\qquad
\cfalse \defeq \Lambda X \ldot \lambda x : X \ldot \lambda y : X \ldot y 
\]
\item Products
\[
\emblue{\tau_1 \times \tau_2 \defeq \forall X \ldot (\tau_1 \to \tau_2 \to X) \to X}
\]
Pairing and projections
\[
\begin{array}{l}
\cpair{e_1, e_2} \defeq \Lambda X \ldot \lambda f : \tau_1 \to \tau_2 \to X \ldot f \, e_1 \, e_2
\\
\cproj_1 \, e \defeq e [\tau_1] \lambda x : \tau_1 \ldot \lambda y : \tau_2 \ldot x
\\
\cproj_2 \, e \defeq e [\tau_2] \lambda x : \tau_1 \ldot \lambda y : \tau_2 \ldot y
\end{array}
\]
\end{itemize}
\end{slide}

\begin{slide}{Representing inductive types}
\begin{itemize}
\item Natural numbers
\[
\emblue{\tnat \defeq \forall X \ldot X \to (X \to X) \to X}
\]
The constants zero and successor
\[
\begin{array}{l}
\czero \defeq \Lambda X \ldot \lambda x : X \ldot \lambda f : X \to X \ldot X , 
\\[0.5em]
\csucc \, n \defeq \Lambda X \ldot \lambda x : X \ldot \lambda f : X \to X \ldot f (n \, X \, x \, f), 
\end{array}
\]
The primitive recursor $\pmb{R}$, of type $\tau \to (\tau \to \tnat \to \tau) \to \tnat \to \tau$ 
can be defined by the iterator $\pmb{It}$ (of type $\tau \to (\tau \to \tau) \to \tnat \to \tau$)
\[
\pmb{It} \ g \ f \ n \defeq n [\tau] g f 
\qquad
\pmb{R} \ g \ f \ n \defeq \cproj_1 (\pmb{It} \, \cpair{g, \czero} \, G \, n) ,
\]
where $G = \lambda x : \tau \times \tnat \ldot \cpair{f (\cproj_1\, x) (\cproj_2 \, x), \csucc (\cproj_2 \, x)}$.
It can be checked that 
\[
\pmb{R} \ g \ f \ \czero \reduceto g , 
\qquad
\pmb{R} \ g \ f \ (\csucc \ n) \reduceto f \ (\pmb{R} \ g \  f  \ n) \ n, 
\]
\end{itemize}
\end{slide}

\begin{slide}{Representing inductive types}
\begin{itemize}
\item List of natural numbers
\[
\emblue{\tlist \defeq \forall X \ldot X \to (\tnat \to X \to X) \to X}
\]
The empty list and constructor
\begin{eqnarray*}
\cnil & \defeq & \Lambda X \ldot \lambda x : X \ldot \lambda f : \tnat \to X \to X \ldot x
\\
\mathtt{cons} \, n \, l & \defeq & \Lambda X \ldot 
\lambda x : X \ldot \lambda f : \tnat \to X \to X \ldot f \, n \, (l \, [X] \, x \, f)
\end{eqnarray*}
A list $[n_1, n_2, \ldots , n_k]$ is basically represented by 
\[
\Lambda X \ldot \lambda x \ldot \lambda f \ldot f \, n_1 ( f \, n_2  \cdots (f \, n_k \, x) \cdots )
\]
Adding a further abstraction we can get a more general version of list (not only of natural numbers) :
\[
\emblue{\forall Y} \ldot \forall X \ldot X \to (\emblue{Y} \to X \to X) \to X 
\]
\end{itemize}
\end{slide}

\begin{slide}{Properties about  system F}
\begin{itemize}
\item System F is strongly normalizing.
\item Not all total recursive functions are representable in system F.
\begin{itemize}
\item The reduction function of system F cannot be programmed in system F.
\end{itemize}
\item According to CH correspondence, system F corresponds to a fragment of \emblue{second-order logic} 
  without the existential quantifier.
\item Type inference in system F is undecidable.
\item In practice, we often use a restricted version of polymorphism where the universal binder only 
  appears at the out-most level of types.
\end{itemize}
\end{slide}

\begin{slide}{Existential types}
\begin{itemize}
\item Existential types supply a mechanism for data abstraction.
\item Existential types
\[
\tau, \tau', \ldots  ::= \ldots \mid \emred{\exists X \ldot \tau}
\]
$\exists X \ldot \tau$ can be seen as an interface of some abstract data type $\tau$, 
whose implementation will be supported by type $X$.
\item Expressions 
\vspace*{-1em}
\[
\begin{array}{lrll}
e, e', \ldots & ::= & \ldots & 
\\
& \mid & \mathtt{pack} \; \sigma \; \mathtt{with} \; e & \mbox{Implementation}
\\
& \mid & \mathtt{open} \; e \; \mathtt{as} \; x \; \mathtt{in} \; e' & 
\mbox{Using abstract type through interface}
\\[-2em]
\end{array}
\]
\item Typing 
\small
\[
\begin{array}{c}
\prooftree
\Delta \vdash \sigma : \Type 
\quad 
\Delta, X : \Type \vdash \tau : \Type
\quad 
\Delta ; \Gamma \vdash e : \tau [\sigma / X]
\justifies
\Delta ; \Gamma \vdash \mathtt{pack} \; \sigma \; \mathtt{with} \; e : \exists X \ldot \tau 
\endprooftree
\\[1.5em]
\prooftree
\Delta ; \Gamma \vdash e : \exists X \ldot \tau 
\quad 
\Delta \vdash \tau' : \Type
\quad 
\Delta, X : \Type ; \Gamma, x: \tau \vdash e' : \tau'
\justifies
\Delta; \Gamma \vdash \mathtt{open} \; e \; \mathtt{as} \; x \; \mathtt{in} \; e' : \tau'
\endprooftree
\end{array}
\]
\end{itemize}
\end{slide}

\begin{slide}{Existential types and data abstraction}
\begin{itemize}
\item Example: an abstract type of stack of natural numbers, consisting of 
\begin{itemize}
\item Formation of the empty stack.
\item Pushing an element at the top of the queue.
\item Popping the top element out.
\end{itemize}
The abstract stack type can be encoded as
\[
\mathsf{Stack} \defeq \exists X \ldot X \times (\tnat \times X \to X) \times (X \to \tnat \times X)
\]
The product defines types of the three operations. \\
An implementation with $\tlist$:
\[
\mathtt{pack} \; \tlist \; \mathtt{with} \;
\cpair{\cnil, \lambda x : \tnat \ldot \lambda y : \tlist  \ldot x :: y, 
\lambda x : \tlist \ldot \cdots }
\vspace*{-2em}
\]
\item Reduction: 
\[
\mathtt{open} \; (\mathtt{pack} \; \sigma \; \mathtt{with} \; e) \; \mathtt{as} \; x \; 
\mathtt{in} \; e' \reduce 
e'[\sigma / X] [e / x]
\]
\end{itemize}
\end{slide}

\begin{slide}{Representing existential types in system F}
\begin{itemize}
\item Existential types do not increase the computational capability of system F: 
they can be encoded in system F.
\begin{itemize}
\item Encoding of types: the $\mathtt{open}$ expression implicitly defines a way of mapping 
  any implementation to whatever types
\[
\exists X \ldot \tau \defeq \forall Y \ldot (\forall X \ldot \tau \to Y) \to Y 
\vspace*{-1.5em}
\]
\item Encoding of expressions:
\begin{eqnarray*}
\mathtt{pack} \; \sigma \; \mathtt{with} \; e & \defeq & 
\Lambda Y \ldot \lambda f : (\forall X \ldot \tau \to Y) \ldot f [\sigma] e 
\\
\mathtt{open} \; e \; \mathtt{as} \; x \; \mathtt{in} \; e' & \defeq & 
e [\tau'] (\lambda x : \tau \ldot e')
\end{eqnarray*}
\end{itemize}
\end{itemize}
\end{slide}

\section{Infinite Data Types}

\begin{slide}{Infinite types}
\begin{itemize}
\item Products and sums are finite types.
\begin{itemize}
\item How do we build infinite types (e.g., $\tnat$) with finite constructors?
\end{itemize}
\item Inductive and co-inductive types are two important forms of infinite types.
\begin{itemize}
\item An {\em inductive type} defines a \emred{least} set of elements that can be produced from 
  a finite set of constructors.
\item A {\em co-inductive type} defines a \emred{largest} set of elements that can be consumed 
  by a finite set of destructors.
\end{itemize}
\end{itemize}
\end{slide}

\begin{slide}{Inductive types}
\begin{itemize}
\item Example: Nat 
\vspace*{-0.5em}
\begin{itemize}
\item Infinitely many element.
\item Every element can be represented by two constructors: $\czero$ and $\csucc$. 
\item The least set containing $\czero$ and closed under $\csucc$. 
\end{itemize}
\item Representing $\tnat$ as an inductive type (by sum)
\[
\prooftree
\Gamma \vdash e : \tunit + \tnat
\justifies
\Gamma \vdash \cfold e : \tnat
\endprooftree
\qquad 
\prooftree
\Gamma \vdash e_1 : \tnat
\quad  
\Gamma, x : \tunit + \tau \vdash e_2 : \tau
\justifies
\Gamma \vdash \mathtt{rec} \; e_1 \; \mathtt{for} \; x \; \mathtt{in} \; e_2 : \tau
\endprooftree
\]
with a reduction rule 
\[
\begin{array}{l}
\mathtt{rec} \; (\emblue{\cfold e_1}) \; \mathtt{for} \; x \; \mathtt{in} \; e_2 
\reduce 
\\ \quad
e_2 [ \mathtt{case}\; e_1 \; \mathtt{of} \; \cinj_1 (\_) \; \mathtt{in} \; \cinj_1 (*) ; 
\\ \hspace*{7em}
\cinj_2 (y) \; \mathtt{in} \;  \cinj_2(\emblue{\mathtt{rec}\; y \;\mathtt{for}\; x \;\mathtt{in}\;  e_2}) / x]
\end{array}
\]
Define $\czero$ and $\csucc$:
$\czero \defeq  \cfold (\cinj_1 *),
\;
\csucc \defeq \lambda x : \tnat \ldot \cfold (\cinj_2 x)$.
\end{itemize}
\end{slide}

\begin{slide}{Co-inductive types}
\begin{itemize}
\item Example: $\tstream$
\vspace*{-0.5em}
\begin{itemize}
\item A stream is an infinite sequence of objects (integers).
\item A stream is characterized by two destructors: 
	$\chd$ returns the head element of the stream and 
	$\ctl$ returns the tail of the stream.
\[
\prooftree
\Gamma \vdash e : \tstream
\justifies
\Gamma \vdash \chd (e) : \tnat
\endprooftree
\qquad \qquad
\prooftree
\Gamma \vdash e : \tstream
\justifies
\Gamma \vdash \ctl (e) : \tstream
\endprooftree
\vspace*{-1em}
\]
\end{itemize}
\item Representing $\tstream$ as a co-inductive type (by product)
\vspace*{-0.5em}
\[
\prooftree
\Gamma \vdash e : \tstream
\justifies
\Gamma \vdash \cunfold  e : \tnat \times \tstream
\endprooftree
\qquad
\prooftree
\Gamma \vdash e_1 : \tau
\quad
\Gamma, x : \tau \vdash e_2 : \tnat \times \tau
\justifies
\Gamma \vdash \mathtt{gen} \; e_1 \; \mathtt{as} \; x \; \mathtt{in} \; e_2 : \tstream
\endprooftree
\]
with a reduction rule
\vspace*{-0.8em}
\[
\begin{array}{l}
\cunfold (\mathtt{gen} \; e_1 \; \mathtt{as} \; x \; \mathtt{in} \; e_2) 
\reduce 
\\ \qquad
\cpair{\cproj_1 (e_2[e_1 / x]), \mathtt{gen} \; \emblue{\cproj_2 (e_2[e_1 /x ])} \; 
  \mathtt{as} \; x \; \mathtt{in} \; e_2}
\\[-0.5em]
\end{array}
\]
Define $\chd$ and $\ctl$: 
$\chd \defeq \lambda x : \tstream \ldot \cproj_1 (\cunfold x), \;
\ctl \defeq \lambda x : \tstream \ldot \cproj_2 (\cunfold x)$.
\end{itemize}
\end{slide}

\begin{slide}{Recursive types}
\begin{itemize}
\item Inductive and co-inductive types can be both seen as \emblue{fix-points} of type operators up to isomorphism.
\begin{itemize}
\item Two types are isomorphic if they are mutually inverse to each other.
\\
We write \emblue{$\mu X \ldot \tau$} for the solution to the equation $X \simeq \tau$.
\item $\tnat$ and $\tunit + \tnat$ are isomorphic and 
	\empurple{$\mu X \ldot \tunit + X$ is a solution to the equation $\tnat \simeq \tunit + \tnat$}.
\item $\tstream$ and $\tnat \times \tstream$ are isomorphic and 
	\empurple{$\mu X \ldot \tnat \times X$ is a solution to the equation $\tstream \simeq \tnat \times \tstream$}.
\item Not every type isomorphism equation has a solution, 
  e.g., $X \simeq X \to \tbool$ (has no solution by Cantor's theorem).
\end{itemize}  
\end{itemize}
\end{slide}

\begin{slide}{Recursive type system}
\begin{itemize}
\item Types ($\lambda^{\to, \mu}$):
\vspace*{-0.5em}
\[
\tau, \tau', \ldots ::= X \mid \emblue{\mu X \ldot \tau} \mid \ldots
\]
$\mu$ is a type binder for recursive types. 
\begin{itemize}
\item Type information judgment $\Delta \vdash \tau : \Type$ 
($\Delta$ is a set of type variables):
\[
\begin{array}{c}
\prooftree
\justifies
\Delta, X: \Type \vdash X : \Type
\endprooftree
\qquad 
\prooftree
\Delta \vdash \tau_1 : \Type
\quad 
\Delta \vdash \tau_2 : \Type
\justifies
\Delta \vdash \tau_1 \to \tau_2 : \Type
\endprooftree
\\[1em]
\ldots \quad \ldots
\\[1em]
\emblue{
\prooftree
\Delta, X : \Type  \vdash \tau : \Type
\justifies
\Delta \vdash \mu X \ldot \tau : \Type
\endprooftree
}
\\[-1em]
\end{array}
\]
\item Typing judgement is of the form $\Delta ; \Gamma \vdash e : \tau$, 
  which indicates \emblue{both $\Delta \vdash \tau : \Type$ and $\Gamma \vdash e : \tau$}
  ($\Gamma$ is the normal typing environment mapping variables to types).
\end{itemize}
\end{itemize}
\end{slide}

\begin{slide}[toc=]{Recursive type system}
\begin{itemize}
\item Expressions
\vspace*{-0.5em}
\[
e ::= \cfold e \mid \cunfold e \mid \ldots
\]
with typing rules:
\[
\prooftree
\Gamma \vdash e : \tau[\mu X . \tau / X]
\justifies 
\Gamma \vdash \cfold e : \tau   
\endprooftree
\qquad 
\prooftree
\Gamma \vdash e : \mu X . \tau
\justifies 
\Gamma \vdash \cunfold e : \tau[\mu X . \tau / X]   
\endprooftree
\]
and reduction rules
\[
\begin{array}{c}
\prooftree
e \reduce e'
\justifies 
\cfold e \reduce \cfold e'
\endprooftree
\qquad 
\prooftree
e \reduce e'
\justifies 
\cunfold e \reduce \cunfold e'
\endprooftree
\qquad
\prooftree
\cfold e \not\reduce 
\justifies 
\cunfold (\cfold e) \reduce e
\endprooftree
\end{array}
\]
The first rule indicates an eager evaluation semantics.
\end{itemize}
\end{slide}

\begin{slide}{Recursive data structure}
\begin{itemize}
\item Natural numbers: $\tnat \defeq \mu X \ldot \tunit + X$.
\[
\czero \defeq \cfold (\cinj_1 *) ,
\quad 
\csucc \defeq \lambda x : \tnat \ldot \cfold (\cinj_2 x)
\]
The conditional $\mathsf{D}$ can be defined using $\mathtt{case}$.
\item List (of natural numbers): $\tlist \defeq \mu X \ldot \tunit + (\tnat \times X)$.
\[
\cnil \defeq \cfold (\cinj_1 *) ,
\quad 
\mathtt{cons} \defeq \lambda x : \tnat \ldot \lambda y : \tlist \ldot \cfold (\cinj_2 \cpair{x, y})
\]
\end{itemize}
\end{slide}

\begin{slide}{Recursion with recursive types}
\begin{itemize}
\item Without further restriction, the system is not strongly normalizing:
\\
Define $\omega = \mu X \ldot X \to X$, then  
\begin{eqnarray*}
& & 
(\lambda x: \omega \ldot (\cunfold x \, x)) (\cfold \lambda x : \omega \ldot (\cunfold x \, x))
\\ & \reduce & 
\cunfold (\cfold \lambda x : \omega \ldot (\cunfold x \, x)) (\cfold \lambda x : \omega \ldot (\cunfold x \, x))
\\ & \reduce & 
(\lambda x : \omega \ldot (\cunfold x \, x)) (\cfold \lambda x : \omega \ldot (\cunfold x \, x))
\\[-2em]
\end{eqnarray*}
Typed $\lambda$-calculus with only \empurple{positive recursive types} is strongly normalizing.
\item Define fixpoint operator (as in PCF) with recursive types: \\
For every type $\tau$, define $\tself(\tau) = \mu X \ldot X \to \tau$.
\vspace*{-1.5em}
\end{itemize}
\small
\begin{eqnarray*}
\cfix_\tau e & \defeq & 
\cunfold (\cfold \lambda x: \tself(\tau) \ldot e (\cunfold x \, x)) 
(\cfold \lambda x: \tself(\tau) \ldot e (\cunfold x \, x)) \\
& \reduce & 
(\lambda x: \tself(\tau) \ldot e (\cunfold x \, x)) 
(\cfold \lambda x: \tself(\tau) \ldot e (\cunfold x \, x)) \\
& \reduce & 
e (\cunfold (\cfold \lambda x: \tself(\tau) \ldot e (\cunfold x \, x)) 
(\cfold \lambda x: \tself(\tau) \ldot e (\cunfold x \, x))) \\
& = & 
e ( \cfix_\tau e)
\end{eqnarray*}
\end{slide}

\end{document}
 

